\documentclass[11pt,a4paper]{article}
\usepackage{amsmath,amssymb,bm,geometry}
\geometry{margin=2.5cm}

\title{Characteristic equations of the clamped--clamped Euler--Bernoulli ant Timoshenko beams}
\author{Stefano Burzio}
\date{}

\begin{document}
\maketitle

This note derives the characteristic equation for free vibration of a Euler-Bernoulli and Timoshenko beam with clamped--clamped boundary conditions. The derivation follows the standard frequency-domain procedure: assume harmonic time dependence, seek exponential spatial solutions, obtain the algebraic characteristic polynomial for the spatial exponent, form the general solution from the roots, impose boundary conditions and state the final determinant condition whose vanishing gives the natural frequencies.

\section{Governing equation of Euler--Bernoulli beams}

Consider a straight prismatic beam of length $\ell$ with Young’s modulus $E$, mass density $\rho$, cross-sectional area $A$, and second moment of area $I$. We denote $w(x,t)$ the transverse displacement of the beam.

In the Euler–Bernoulli beam theory, cross-sections remain perpendicular to the neutral axis during bending, so the rotation of the cross-section is equal to the slope of the deflection, $\varphi(x,t) = w’(x,t)$. Neglecting shear deformation and rotary inertia, the linear equation of motion for free vibration is
\[
      EI\,\frac{d^4 w(x,t)}{dx^4} = \rho A\,\frac{d^2 w(x,t)}{dt^2}.
\]
Assuming harmonic motion
\[
      w(x,t) = W(x)e^{i\omega t},
\]
the governing differential equation becomes
\[
      EI\,W''''(x) = \rho A \omega^2 W(x).
\]
Define the wavenumber parameter
\[
      \beta^{4} = \frac{\rho A }{EI}\omega^2,
\]
so that
\[
      W''''(x) - \beta^{4} W(x) = 0.
\]

\subsection{General solution of Euler--Bernoulli equation}

The general solution of this constant--coefficient differential equation is
\[
      W(x) = A_1\cosh(\beta x) + A_2\sinh(\beta x)
      + A_3\cos(\beta x) + A_4\sin(\beta x).
\]

\subsection{Boundary conditions}

For a clamped end:
\[
      w(0)=w(\ell)= 0, \qquad \frac{dw}{dx} (0)=\frac{dw}{dx} (\ell)= 0.
\]
Thus, for a clamped--clamped beam:
\[
      W(0) = W'(0) = W(\ell) = W'(\ell) = 0.
\]
At $x=0$:
\[
      \begin{aligned}
            W(0)  & = A_1 + A_3 = 0,        \\
            W'(0) & = \beta(A_2 - A_4) = 0.
      \end{aligned}
\]
Hence
\[
      A_3 = -A_1, \qquad A_4 = A_2.
\]
Substituting these into $W(x)$ gives
\[
      W(x) = A_1(\cosh\beta x - \cos\beta x)
      + A_2(\sinh\beta x - \sin\beta x).
\]
At $x=\ell$:
\[
      \begin{aligned}
            W(\ell)  & = A_1(\cosh\beta\ell- \cos\beta \ell)
            + A_2(\sinh\beta\ell- \sin\beta \ell) = 0,              \\[4pt]
            W'(\ell) & = \beta\Big[A_1(\sinh\beta\ell+ \sin\beta L)
                  + A_2(\cosh\beta\ell- \cos\beta L)\Big] = 0.
      \end{aligned}
\]
These two equations can be written in matrix form:
\[
      \begin{bmatrix}
            \cosh\beta\ell- \cos\beta\ell & \sinh\beta\ell- \sin\beta\ell \\
            \sinh\beta\ell+ \sin\beta\ell & \cosh\beta\ell- \cos\beta L
      \end{bmatrix}
      \begin{bmatrix}
            A_1 \\[3pt] A_2
      \end{bmatrix}
      = \bm{0}.
\]

\subsection{Characteristic equation for Euler-Bernoulli beams}
For a nontrivial solution $(A_1, A_2) \neq (0,0)$, the determinant of the
coefficient matrix must vanish:
\[
      \begin{vmatrix}
            \cosh\beta\ell- \cos\beta\ell & \sinh\beta\ell- \sin\beta\ell \\
            \sinh\beta\ell+ \sin\beta\ell & \cosh\beta\ell- \cos\beta L
      \end{vmatrix} = 0.
\]

Expanding the determinant gives
\[
      (\cosh\beta\ell- \cos\beta L)^2
      - (\sinh\beta\ell+ \sin\beta L)(\sinh\beta\ell- \sin\beta L) = 0.
\]

Simplifying using trigonometric and hyperbolic identities,
\[
      \boxed{\cosh(\beta L)\cos(\beta L) - 1 = 0.}
\]

\subsection{Natural frequencies for Euler-Bernoulli beams}
Let $\beta_n L$ be the roots of the above transcendental equation.
Typical values are
\[
      \beta_1\ell= 4.7300, \quad
      \beta_2\ell= 7.8532, \quad
      \beta_3\ell= 10.9956, \quad
      \beta_4\ell= 14.1372, \ \ldots
\]

The corresponding circular natural frequencies are
\[
      \omega_n = \beta_n^2 \sqrt{\frac{EI}{\rho A \ell^4}},
\]
and the natural frequencies in hertz are
\[
      f_n = \frac{\omega_n}{2\pi}.
\]
The mode shapes are given by
\[
      W_n(x) = A \big[\cosh(\beta_n x) - \cos(\beta_n x)\big]
      - B \big[\sinh(\beta_n x) - \sin(\beta_n x)\big],
\]
with constants $A$ and $B$ determined up to a normalization factor
from the boundary conditions at $x=\ell$.

\section{Governing equations for Timoshenko beams}

Consider a straight prismatic beam of length \(l\) with Young's modulus \(E\), shear modulus \(G\), mass density \(\rho\), cross-sectional area \(A\), second moment of area \(I\), and shear correctionfactor \(k\). We denote \(w(x,t)\) the transverse displacement and \(\varphi(x,t)\) the cross-section rotation (slope of the cross-section, not necessarily equal to \(w'\)).

The linear Timoshenko beam theory (including shear deformation and rotary inertia)
gives the following coupled equations of motion (signs depend on convention).
Using the convention below, in the time domain the equations read
$$
      \begin{aligned}
            \frac{\partial}{\partial x}\big( kGA(\varphi - w')\big) + \rho A\,\ddot w                    & = 0,\
            \frac{\partial}{\partial x}\big(EI\,\varphi'\big) + kGA(\varphi - w') + \rho I\,\ddot\varphi & = 0,
      \end{aligned}
$$
where \( '\equiv\partial/\partial x\) and \(\ddot{}\equiv\partial^2/\partial t^2\),
and \(\rho I\) is the rotary inertia per unit length.

Eliminating $w$ or $\varphi$ from the previous equations, the following two uncoupled differential equations in $w$ and $\varphi$ are obtained:
$$
      \begin{aligned}
             & E I \frac{\partial^4 w}{\partial x^4} + \rho A \frac{\partial^2 w}{\partial t^2} - \left(\rho I + \frac{E I}{k G}\right) \frac{\partial^4 w}{\partial x^2 \partial t^2} + \frac{\rho I}{k G} \frac{\partial^4 w}{\partial t^4} = 0,                        \\[1em]
             & E I \frac{\partial^4 \varphi}{\partial x} + \rho A \frac{\partial^2 \varphi}{\partial t^2} - \left(\rho I + \frac{E I}{k G} \right) \frac{\partial^4 \varphi}{\partial x^2 \partial t^2} + \frac{\rho I}{k G} \frac{\partial^4 \varphi}{\partial t^4} = 0.
      \end{aligned}
$$
Assume harmonic time dependence :
\[
      w(x,t)=W(\xi)\,e^{i\omega t},\qquad \varphi(x,t)=\Phi(\xi)\,e^{i\omega t},
\]
with $\xi=x/\ell$ and substitute into the uncoupled equations. This yields the frequency-domain system
$$
      \begin{aligned}
             & W^{\mathrm{iv}}(\xi) + b^2\left(r^2 + s^2\right) W^{\prime \prime}(\xi) - b^2 \left(1 - b^2 r^2 s^2\right) W(\xi) = 0,          \\[1em]
             & \Phi^{\mathrm{iv}}(\xi) + b^2\left(r^2 + s^2\right) \Phi^{\prime \prime}(\xi) - b^2 \left(1 - b^2 r^2 s^2\right) \Phi(\xi) = 0.
      \end{aligned}
$$
where

\begin{align}
      b^2 & =\frac{\rho A}{E I} \ell ^4 \omega^2, \label{btoomega} \\
      r^2 & =\frac{I}{A \ell^2}, \nonumber                         \\
      s^2 & =\frac{E I}{k A G \ell^2}.  \nonumber
\end{align}

\subsection{General solution of Timoshenko equation}

The general solution of this two constant--coefficient fourth order differential equations are
\begin{align}
      \label{genTimo1}
      W(\xi)    & =C_1 \cosh b \alpha \xi+C_2 \sinh b \alpha \xi+C_3 \cos b \beta \xi +C_4 \sin b \beta \xi,                                              \\
      \Psi(\xi) & =C_{1}^{\prime} \sinh b \alpha \xi+C_{2}^{\prime} \cosh b \alpha \xi+C_{3 }^{\prime} \sin b \beta \xi +C_{4}^{\prime} \cos b \beta \xi,
      \label{genTimo2}
\end{align}
where
$$
      \begin{aligned}
            \alpha & =\frac{1}{\sqrt{ 2}}\left\{-\left(r^2+s^2\right)+\left[\left(r^2-s^2\right)^2+4 / b^2\right]^{1 / 2}\right\}^{1 / 2} \\
            \beta  & =\frac{1}{\sqrt{ 2} }\left\{\left(r^2+s^2\right)+\left[\left(r^2-s^2\right)^2+4 / b^2\right]^{1 / 2}\right\}^{1 / 2}
      \end{aligned}
$$
and
$$
      \left[\left(r^2-s^2\right)^2+4 / b^2\right]^{1 / 2}>\left(r^2+s^2\right)
$$
is assumed.

\subsection{Boundary conditions}
For a clamped end:
\[
      w(0)=w(\ell)= 0, \qquad \varphi (0)=\varphi(\ell)= 0.
\]
Thus, for a clamped--clamped beam:
\[
      W(0) = \Psi(1) = W(\ell) = \Psi(1) = 0.
\]

\subsection{Characteristic equation for Timoshenko beams}
The application of appropriate boundary conditions to equations \eqref{genTimo1} and \eqref{genTimo2} yields to a set of four homogeneous linear algebraic equations in four constants $C_{1}$ to $C_{4}$  and $C_{1}^\prime$ to $C_{4}^\prime$. In order that the solutions other than zero may exist the determinant of the coefficients of $C$, must be equal to zero. This leads to the frequency equation in each case from which the natural frequencies can be determined:
$$
      \boxed{2-2 \cosh b \alpha \cos b \beta +\frac{b}{\left(1-b^2 r^2 s^2\right)^{1 / 2}}
            \left[b^2 s^2\left(r^2-s^2\right)^2+\left(3 s^2-r^2\right)\right] \sinh b \alpha \sin b \beta=0}
$$
This characteristic equation admits a sequence of positive solutions $\beta_n$. To compute the natural frequencies, we need to invert \eqref{btoomega} to obtain the natural frequencies $\omega_n$:
\[
      \omega_n = \frac{\beta_i}{\ell^2} \sqrt{\frac{EI}{\rho A}} \qquad n=1,2,\ldots.
\]

\bigskip
\subsection*{References.}
Huang (1961), "\emph{Transverse Vibration of Beams Including Rotary Inertia and Shear Deformation}."

\end{document}